\documentclass[10pt,a4paper]{amsart}

\pdfinfoomitdate=1
\pdftrailerid{}
\pdfsuppressptexinfo=15

\usepackage[T1]{fontenc}
\usepackage{lmodern}
\usepackage[margin=25mm]{geometry}
\usepackage[protrusion=true,expansion=false]{microtype}
\usepackage{amsmath,amssymb,mathtools}
\usepackage{booktabs,array}
\usepackage{enumitem}
\usepackage{float}
\usepackage{xcolor}
\usepackage[numbers,sort&compress]{natbib}
\usepackage[colorlinks=true,linkcolor=blue!55!black,citecolor=blue!55!black,urlcolor=blue!55!black]{hyperref}
\usepackage[nameinlink,noabbrev]{cleveref}

\setlist{nosep,leftmargin=*}
\raggedbottom
\allowdisplaybreaks
\emergencystretch=2em

\newtheorem{theorem}{Theorem}[section]
\newtheorem{proposition}[theorem]{Proposition}
\newtheorem{lemma}[theorem]{Lemma}
\newtheorem{corollary}[theorem]{Corollary}
\theoremstyle{definition}
\newtheorem{definition}[theorem]{Definition}
\theoremstyle{remark}
\newtheorem{remark}[theorem]{Remark}

\newcommand{\E}{\mathbb E}
\newcommand{\Pp}{\mathbb P}
\newcommand{\supp}{\operatorname{supp}}
\newcommand{\one}{\mathbf 1}

\title[Higher-moment poles in coalescence]{Higher-Moment Pole Ladders in Product-Plus-One Coalescence:\\A Yule-Averaged Marked Refinement}
\author{Anonymous}
\date{}
\subjclass[2020]{60C05, 05C05, 05A15, 05A16}
\keywords{random coalescence, binary search tree, antichain polynomial,
Riccati equation, moment generating function, singularity analysis}

\begin{document}

\begin{abstract}
Start with $n$ ordered copies of $1$, repeatedly choose a current adjacent
pair uniformly, and replace $(x,y)$ by $xy+1$; write $X_n$ for the terminal
value.  Disanto, Fuchs, Paningbatan, and Rosenberg proved that the number
$R_n$ of root ancestral configurations in a Yule--Harding random tree obeys
the same recursion after the shift $X_n=R_n+1$.  We prove the identification
under the common ordered-history coupling and assign their split law, finite
distribution, unmarked antichain interpretation, mean analysis, and
second-moment neighborhood zero contribution credit.  After this
subtraction, and after crediting the existing fixed-tree cardinality marker,
we obtain its Yule-averaged bivariate transform in closed form.  The
all-order moment equation is a mechanical expansion of the owned exact law.
The main residual theorem continues the two owned low-order radii to a strict ladder:
for every $r\geq3$, the $r$th raw-moment series has a normalized positive
simple pole at a radius $\rho_r<\rho_{r-1}$ and
$\limsup_n(\E X_n^r)^{1/n}=\rho_r^{-1}$.  No full coefficient asymptotic is
claimed for $r\geq3$.  The residual has passed only a bounded owner
audit; novelty, priority, and external release remain open.
\end{abstract}

\maketitle

\section{Introduction and scope}

Consider the following rank-decreasing random dynamics.  Begin with the
ordered word $(1,\ldots,1)$ of length $n$.  If the current word has $j$
entries, choose one of its $j-1$ adjacent pairs with probability $1/(j-1)$
and apply
\begin{equation}\label{eq:merge}
                 (x,y)\longmapsto xy+1.
\end{equation}
After $n-1$ mergers a positive integer $X_n$ remains.  The local rule is
nonassociative, so $X_n$ depends on the entire merger genealogy.  The first
nondegenerate law already occurs at $n=4$:
\[
        \Pp(X_4=4)=\frac23,\qquad \Pp(X_4=5)=\frac13.
\]

The decisive owner is not merely a generic random-tree neighbor.
\citet{DisantoEtAl2022} study the number $R_n$ of root ancestral
configurations under the Yule--Harding model.  Their Proposition~3.5 gives
$R_1=0$ and the uniform-split recurrence
$R_n\overset d=(R_I+1)(R'_{n-I}+1)$, including the exact finite law.  Their
Section~3.2 identifies root configurations with nonempty antichains of the
pruned internal-node tree; Sections~5.2--5.3 derive the mean and
second-moment Riccati analyses.  Thus $X_n=R_n+1$ is the same statistic in
shifted notation.  All of those statements, including the mean pole
$1-\exp(-2\pi/(3\sqrt3))$ and the order-two growth analysis, receive zero
contribution credit here.

Two further owners narrow the residual.  The cardinality refinement on a
fixed rooted-tree poset is studied by
\citet{AndriantianaWagnerWang2013}; only its average under the Yule/uniform
ordered-history law and the resulting closed bivariate transform are at
issue here.  Moreover, \citet[Table~1]{ChangFuchs2010} give the exact
Yule--Harding probability $2^{n-2}/(n-1)!$ of an $n$-caterpillar, building on
the caterpillar analysis of \citet{Rosenberg2006}.  Both the minimizing shape
and its probability therefore receive zero credit.

The genealogy itself is also classical.  Each current boundary is an
undeleted original boundary, and successive uniform choices form a uniform
permutation.  Its max-Cartesian tree has the usual random binary-search-tree
(BST) root-split law.  The same subtraction applies to deterministic subtree enumeration
\cite{Ruskey1981Subtrees}, ideals of forest posets
\cite{KodaRuskey1993Ideals}, antichain statistics on rooted plane trees
\cite{Klazar1997Twelve}, and random forest ideals
\cite{Janson2002Ideals}.  Random-BST generating-function schemes and
singularity methods are also established
\cite{FlajoletGourdonMartinez1997Patterns,MartinezPanholzerProdinger1998Descendants,FlajoletSedgewick2009}.

After those deductions, the paper records only the following residual
package:
\begin{enumerate}[label=\textup{(\roman*)}]
\item the Yule-averaged bivariate ordinary generating function for the
      owned fixed-tree cardinality marker, together with its closed Euler
      form; and
\item beginning at $r=3$, a strict continuation of the moment-radius ladder
      with unit residues and exact exponential limsups.  The all-order
      triangular identity is retained only as the mechanical interface from
      the owned exact law.
\end{enumerate}
The higher-moment conclusion is deliberately limited to a positive simple
pole and an exact exponential \emph{limsup}.  Without control of all complex
singularities, no full asymptotic is asserted for $r\geq3$.  The direct owner
has a stronger order-two conclusion than the ceiling used here.

Two complementary viewpoints organize the proofs.  The temporal viewpoint
uses uniform deletion orders, conditional root splits, moments, and
Riccati--Euler equations.  The structural viewpoint regards the same orders
as heap labelings of Cartesian trees and interprets the output as a
cardinality-marked antichain count.  Both viewpoints use the same split law;
we do not present them as probabilistically independent proofs.  The second
viewpoint supplies a Yule-averaged cardinality transform and exact comb
controls rather than a new random-tree model.  A focused search of the direct owner's
citation neighborhood and current root-configuration formulations did not
locate the averaged transform or a strict $r\geq3$ continuation of the owned
low-order pole ladder.  This bounded non-hit
is not evidence of novelty or priority.
External circulation remains \textsc{hold}.

\section{Direct-owner interface and the exact shift}

Label the $n-1$ original boundaries by $1,\ldots,n-1$.  A merger deletes
exactly the boundary separating the two merged blocks; every other current
boundary retains its original label.

\begin{theorem}[Owned split law and objectwise shift]\label{thm:split}
Let $I_n$ be uniform on $\{1,\ldots,n-1\}$.  Conditional on $I_n=i$, let
$X_i$ and $X'_{n-i}$ be independent copies generated on the two sides of the
last-deleted boundary.  Then
\begin{equation}\label{eq:split}
 X_1=1,\qquad
 X_n\ \overset d=\ 1+X_{I_n}X'_{n-I_n}\quad(n\geq2).
\end{equation}
Consequently, if $p_n(v)=\Pp(X_n=v)$, then $p_1(v)=\one_{\{v=1\}}$ and
\begin{equation}\label{eq:law}
 p_n(v)=\frac1{n-1}\sum_{i=1}^{n-1}
 \sum_{\substack{a,b\geq1\\1+ab=v}}p_i(a)p_{n-i}(b).
\end{equation}
Every law in \eqref{eq:law} has finite support.
Moreover, couple the merger process to the ordered unlabeled history having
the same boundary order.  If $R_n$ counts its root configurations, then
\begin{equation}\label{eq:shift}
                         X_n=R_n+1
\end{equation}
for every history, not only in distribution.
\end{theorem}

\begin{proof}
When $j$ blocks remain, precisely $j-1$ original boundaries remain, and the
process selects each with probability $1/(j-1)$.  Thus the complete deletion
order is a uniform permutation of $\{1,\ldots,n-1\}$.  Its last element is
uniform.  Conditional on that element being $i$, the relative orders of the
$i-1$ left boundaries and the $n-i-1$ right boundaries are independent and
uniform.  All mergers on each side occur before the last deletion, after
which \eqref{eq:merge} combines the two terminal side values.  This proves
\eqref{eq:split}.  Conditioning on $I_n$ and the two side values gives
\eqref{eq:law}.  Finiteness follows inductively.  These distributional
statements reproduce Proposition~3.5 of \citet{DisantoEtAl2022} after a
shift and receive zero credit.

For the objectwise identification, let $T$ be the common ordered evaluation
tree.  At a leaf, its root-configuration count is $R(T)=0$, while the merger
value is $V(T)=1$.  At an internal root, the standard root-configuration
decomposition gives
\[
 R(T)=\bigl(R(T_L)+1\bigr)\bigl(R(T_R)+1\bigr).
\]
Induction and $V(T)=1+V(T_L)V(T_R)$ now give $V(T)=R(T)+1$.  This proves
\eqref{eq:shift}; the root-configuration recurrence and its unmarked
antichain interpretation are owned by \citet{DisantoEtAl2022}.
\end{proof}

The proof also fixes the orientation: the \emph{last} deleted boundary is the
root of the evaluation tree.  Reversing the deletion order exchanges maxima
and minima in the Cartesian-tree description and is not the process above.

\section{Heap labelings and a marked antichain transform}

For a fixed deletion order, let $T$ be the ordered full binary evaluation
tree.  Its leaves carry $1$, and every internal vertex applies
$V(T)=1+V(T_L)V(T_R)$.  Order the internal vertices by ancestry.  For an
indeterminate $s$, define
\[
 P_T(s)=\sum_{B} s^{|B|},
\]
where $B$ ranges over antichains of internal vertices, including the empty
antichain.  Put $P_{\mathrm{leaf}}(s)=1$.

\begin{lemma}[Owned unmarked interface and its marker]\label{lem:antichain}
For every evaluation tree,
\begin{equation}\label{eq:treepoly}
 P_T(s)=s+P_{T_L}(s)P_{T_R}(s),\qquad V(T)=P_T(1).
\end{equation}
\end{lemma}

\begin{proof}
An antichain either contains the root, in which case it is the singleton
root, or excludes the root, in which case its restrictions to the two
subtrees are arbitrary antichains and can be chosen independently.  The
first identity follows.  Its specialization at $s=1$ is the same recursion
and initial condition as $V(T)$.
\end{proof}

At $s=1$, this is precisely the direct owner's correspondence: $R(T)$
counts nonempty antichains and $X(T)=R(T)+1$ also includes the empty one.
For general $s$, the fixed-tree cardinality marker and the equivalent
root-subtree leaf statistic are also owned by
\citet{AndriantianaWagnerWang2013}.  Only averaging this marker under the
Yule/uniform ordered-history law and solving the resulting bivariate equation
belong to the residual claim set.

Let $T_n$ denote the tree induced by a uniform boundary order and write
\[
 a_n(s)=\E P_{T_n}(s),\qquad
 A(z,s)=\sum_{n\geq1}a_n(s)z^{n-1}.
\]

\begin{theorem}[Expected cardinality-marked antichains]\label{thm:marked}
Coefficientwise as a formal series in $z$ and $s$,
\begin{equation}\label{eq:markedriccati}
 \partial_z A=A^2+\frac{s}{(1-z)^2},\qquad A(0,s)=1.
\end{equation}
Set $w=1-z$,
\[
 \delta=\sqrt{1-4s},\qquad
 \beta_\pm=\frac{1\pm\delta}{2},
\]
and
\begin{equation}\label{eq:Y}
 Y(w,s)=\frac{\beta_+w^{\beta_+}-\beta_-w^{\beta_-}}{\delta}.
\end{equation}
With the removable value
\[
 Y(w,1/4)=w^{1/2}\left(1+\frac12\log w\right),
\]
the solution is
\begin{equation}\label{eq:markedclosed}
                     A(z,s)=\frac{Y_w(1-z,s)}{Y(1-z,s)}.
\end{equation}
In particular, $[z^{n-1}s^k]A$ is the exact expected number of
$k$-element internal-node antichains, and $A(z,1)$ is the mean series of
$X_n$.
\end{theorem}

\begin{proof}
The split law and \cref{lem:antichain} give, for $n\geq2$,
\begin{equation}\label{eq:markedrec}
 a_n(s)=s+\frac1{n-1}\sum_{i=1}^{n-1}a_i(s)a_{n-i}(s).
\end{equation}
Multiplying by $(n-1)z^{n-2}$ and summing gives
\eqref{eq:markedriccati}.  The function in \eqref{eq:Y} satisfies
\[
 Y_{ww}+\frac{s}{w^2}Y=0,\qquad Y(1,s)=Y_w(1,s)=1.
\]
Because $\partial_z=-\partial_w$, the logarithmic derivative
$A=Y_w/Y$ satisfies \eqref{eq:markedriccati} and the initial condition.
Uniqueness of formal solutions proves \eqref{eq:markedclosed}.  Taking the
limit $\delta\to0$ in \eqref{eq:Y} gives the displayed removable value.
The coefficient interpretation is the definition of $a_n(s)$, and
\cref{lem:antichain} gives $a_n(1)=\E X_n$.
\end{proof}

Equation \eqref{eq:markedclosed} is interpreted formally near $s=0$, or
analytically with consistent branches.  At $s=0$ it gives
$A(z,0)=(1-z)^{-1}$, recording the unique empty antichain.  The marked
transform adds a coefficient-level statement, but its derivation still uses
the same classical root split as \cref{thm:split}.

\section{All raw moments and a strict radius cascade}

For integers $r\geq0$ set
\[
 m_{r,n}=\E X_n^r,\qquad
 F_r(z)=\sum_{n\geq1}m_{r,n}z^{n-1}.
\]
Thus $m_{0,n}=1$ and $F_0(z)=(1-z)^{-1}$.

\begin{proposition}[Mechanical all-order moment interface]\label{prop:moments}
For $n\geq2$ and $r\geq0$,
\begin{equation}\label{eq:momentrec}
 (n-1)m_{r,n}=\sum_{i=1}^{n-1}\sum_{k=0}^r
       \binom rk m_{k,i}m_{k,n-i}.
\end{equation}
Equivalently,
\begin{equation}\label{eq:momentode}
 F_r'(z)=\sum_{k=0}^r\binom rk F_k(z)^2,
 \qquad F_r(0)=1.
\end{equation}
At level $r$, the only new function on the right is $F_r^2$.
\end{proposition}

\begin{proof}
Condition on the split in \eqref{eq:split}, expand
$(1+ab)^r=\sum_{k=0}^r\binom rk a^k b^k$, and use conditional
independence.  This gives \eqref{eq:momentrec}.  Coefficient extraction from
the square of each $F_k$ gives \eqref{eq:momentode}.
\end{proof}

The cases $r=1,2$ lie in the mean/variance analysis of
\citet{DisantoEtAl2022}.  More generally, the displayed identity is a
one-line binomial expansion of their exact distributional recurrence.  It
therefore receives zero contribution credit and serves only as the interface
for the strict higher-order theorem; no elementary solution is claimed at
every level.

\begin{theorem}[Strict higher-moment radius continuation]\label{thm:radii}
There are numbers
\begin{equation}\label{eq:rholadder}
             1=\rho_0>\rho_1>\rho_2>\cdots>0
\end{equation}
such that $\rho_r$ is the radius of convergence of $F_r$.  For every
$r\geq1$,
\begin{equation}\label{eq:localpole}
 F_r(z)=\frac1{\rho_r-z}+O(1)\qquad(z\to\rho_r),
\end{equation}
and
\begin{equation}\label{eq:limsup}
 \limsup_{n\to\infty}\bigl(\E X_n^r\bigr)^{1/n}=\rho_r^{-1}.
\end{equation}
The residual assertion begins at $r=3$: the cases $r=1,2$ are retained only
as the owned base of the continuation.
\end{theorem}

\begin{proof}
Proceed by induction from $F_0(z)=(1-z)^{-1}$ and $\rho_0=1$.  For
$r\geq1$, separate the new unknown in \eqref{eq:momentode}:
\begin{equation}\label{eq:Gr}
 F_r'=F_r^2+G_r,\qquad
 G_r=\sum_{k=0}^{r-1}\binom rk F_k^2.
\end{equation}
On $|z|<\rho_{r-1}$, let $U_r$ solve
\begin{equation}\label{eq:Ur}
 U_r''+G_rU_r=0,\qquad U_r(0)=1,\quad U_r'(0)=-1.
\end{equation}
The logarithmic derivative $-U_r'/U_r$ satisfies \eqref{eq:Gr} and agrees
with $F_r$ near zero.

By induction, all terms in $G_r$ except the one containing $F_{r-1}$ are
analytic at $\rho_{r-1}$, while \eqref{eq:localpole} gives
\begin{equation}\label{eq:Gedge}
 G_r(x)=\frac{r}{(\rho_{r-1}-x)^2}
        +O\!\left(\frac1{\rho_{r-1}-x}\right)
 \quad (x\uparrow\rho_{r-1}).
\end{equation}
Choose $c$ with $1/4<c<r$.  Close enough to $\rho_{r-1}$,
$G_r(x)\geq c(\rho_{r-1}-x)^{-2}$.  Every nonzero solution of the Euler
comparison equation
\[
 V''+\frac{c}{(\rho_{r-1}-x)^2}V=0
\]
has infinitely many zeros accumulating at $\rho_{r-1}$: after putting
$s=\rho_{r-1}-x$, it is a linear combination of
$s^{1/2}\cos(\nu\log s)$ and $s^{1/2}\sin(\nu\log s)$, with
$\nu=\sqrt{4c-1}/2$.  If $U_r$ has not already vanished, choose two
consecutive comparison zeros after the point where the coefficient
inequality holds.  The Sturm comparison theorem forces a zero of $U_r$
strictly between them \cite[Chapter~1]{Zettl2005}.  Hence $U_r$ vanishes
strictly before $\rho_{r-1}$.  Let $\rho_r$ be its first positive zero.
Then $0<\rho_r<\rho_{r-1}$.

The coefficient $G_r$ is analytic at $\rho_r$.  A double zero of $U_r$
would contradict uniqueness for \eqref{eq:Ur}, so the zero is simple.
Consequently $-U_r'/U_r$ has the normalized local form
\eqref{eq:localpole}.

It remains to identify the positive pole with the complex radius.  The
coefficients of $F_r$ are positive.  If its radius $R$ were smaller than
$\rho_r$, Pringsheim's theorem would force a singularity at the positive
point $R$ \cite[Section~IV.6]{FlajoletSedgewick2009}.  Yet $G_r$ is analytic
there and $U_r$ has no positive zero before $\rho_r$, so $-U_r'/U_r$ is
analytic in a neighborhood of $R$, a contradiction.  Thus $R=\rho_r$.
Cauchy--Hadamard now gives \eqref{eq:limsup}; replacing the exponent
$1/(n-1)$ by $1/n$ does not change the limit superior.
\end{proof}

\begin{remark}[Hard ceiling for $r\geq3$]\label{rem:ceiling}
The local pole in \eqref{eq:localpole} lies at the positive convergence
radius, but the proof does not exclude other singularities of the same
modulus.  Accordingly, \cref{thm:radii} claims only \eqref{eq:limsup} for
$r\geq3$.  A statement such as
$\E X_n^r\sim\rho_r^{-n}$ would require an additional dominant-singularity
theorem and is not part of this paper.  The stronger order-two asymptotic in
\citet{DisantoEtAl2022} is not weakened or reclaimed.
\end{remark}

\section{Owned mean and endpoint normalizations}

At $r=1$, the forcing in \eqref{eq:Gr} is explicit.  This permits a complete
coefficient asymptotic, unlike the higher levels.

\begin{theorem}[Direct-owner mean, reproduced for normalization]\label{thm:mean}
Let $M=F_1$ and $w=1-z$.  Then
\begin{equation}\label{eq:meanode}
 M'=M^2+\frac1{(1-z)^2},\qquad M(0)=1,
\end{equation}
and
\begin{equation}\label{eq:meanclosed}
 M(z)=\frac1w\left[
 \frac12-\frac{\sqrt3}{2}
 \tan\!\left(\frac{\sqrt3}{2}\log w-\frac\pi6\right)
 \right].
\end{equation}
Its unique dominant singularity is the simple pole
\begin{equation}\label{eq:rho}
 \rho=1-\exp\!\left(-\frac{2\pi}{3\sqrt3}\right)
      =0.7015639408\ldots .
\end{equation}
For some $R>\rho$,
\begin{equation}\label{eq:meanasymptotic}
 \E X_n=\rho^{-n}+O(R^{-n}),\qquad
 \rho^{-1}=1.4253868276\ldots .
\end{equation}
In particular, the leading coefficient is one.
\end{theorem}

\begin{proof}
Equation \eqref{eq:meanode} is the case $r=1$ of
\cref{prop:moments}.  Write $M=-U'/U$.  Then
\[
 U''+\frac{U}{(1-z)^2}=0,\qquad U(0)=1,\quad U'(0)=-1.
\]
With $\alpha=\sqrt3/2$, the solution is
\begin{equation}\label{eq:Umean}
 U(z)=\frac2{\sqrt3}w^{1/2}
 \cos\!\left(\alpha\log w-\frac\pi6\right).
\end{equation}
Since $M=U_w/U$, differentiation yields \eqref{eq:meanclosed}.

In $|z|<1$, the number $w=1-z$ lies in the right half-plane, so the
principal logarithm is analytic.  The zeros of the cosine in
\eqref{eq:Umean} satisfy
\[
 w_k=\exp\!\left[\frac2{\sqrt3}
       \left(\frac{2\pi}{3}+k\pi\right)\right],
 \qquad k\in\mathbb Z.
\]
They are positive real because all complex zeros of cosine are real.  The
zero $k=-1$ gives \eqref{eq:rho}.  The zeros with $k\leq-2$ correspond to
larger positive $z<1$, while those with $k\geq0$ give $|z|>1$; the logarithm
first branches at $z=1$.  Hence $\rho$ is the unique singularity on its
circle and is isolated from the next one.  The logarithmic derivative of a
simple zero has local form $1/(\rho-z)+O(1)$.  Coefficient extraction gives
\eqref{eq:meanasymptotic}.
\end{proof}

The theorem is the shifted form of the Yule--Harding mean analysis in
\citet[Section~5.2]{DisantoEtAl2022}; no part of it is counted as a residual
advance.

The opposite edge of the distribution also has a closed form.

\begin{proposition}[Fully owned caterpillar normalization]\label{prop:minimum}
For every $n\geq1$, $X_n\geq n$.  For $n\geq2$,
\begin{equation}\label{eq:minatom}
       \min\supp X_n=n,
       \qquad
       \Pp(X_n=n)=\frac{2^{n-2}}{(n-1)!}.
\end{equation}
\end{proposition}

\begin{proof}
For positive integers, $xy+1\geq x+y$, with equality exactly when $x=1$ or
$y=1$.  Induct on $n$ in \eqref{eq:split}.  If the root split has sizes
$i$ and $j=n-i$, then
\[
 1+X_iX'_j\geq1+ij=i+j+(i-1)(j-1)\geq n.
\]
Equality requires minimal values in both subtrees and an endpoint split
$i=1$ or $j=1$.  Iterating this condition gives precisely the planar combs.
Among the $(n-1)!$ boundary orders, their number $c_n$ satisfies
$c_2=1$ and $c_n=2c_{n-1}$ for $n\geq3$.  Thus
$c_n=2^{n-2}$, which proves \eqref{eq:minatom}.
\end{proof}

The fact that the caterpillar minimizes the number of root configurations is
already recorded by \citet[Section~2.4.4]{DisantoEtAl2022}.  The displayed
mass is exactly the Yule--Harding $n$-caterpillar probability printed by
\citet[Table~1]{ChangFuchs2010}; \citet{Rosenberg2006} is the earlier
caterpillar-pattern owner.  The proof is retained only as a translation to
the adjacent-deletion encoding, and the entire proposition receives zero
contribution credit.

\section{Exact coefficient controls and proof dependencies}

The first marked polynomials make the coefficient interpretation concrete.
They were obtained both from \eqref{eq:markedrec} and by averaging over every
boundary order through $n=9$.

\begin{table}[H]
\centering
\small
\caption{Exact expected cardinality-marked antichain polynomials.  At
$s=1$, the owned root-configuration statistic satisfies $R_n=X_n-1$.
The machine-readable artifact continues through $n=12$.}
\label{tab:marked}
\begin{tabular}{@{}rll@{}}
\toprule
$n$ & $a_n(s)=\E P_{T_n}(s)$ & $\E X_n$ \\
\midrule
1 & $1$ & $1$ \\
2 & $1+s$ & $2$ \\
3 & $1+2s$ & $3$ \\
4 & $1+3s+\frac13s^2$ & $\frac{13}{3}$ \\
5 & $1+4s+\frac76s^2$ & $\frac{37}{6}$ \\
6 & $1+5s+\frac{13}{5}s^2+\frac{2}{15}s^3$ & $\frac{131}{15}$ \\
7 & $1+6s+\frac{47}{10}s^2+\frac{59}{90}s^3$ & $\frac{556}{45}$ \\
8 & $1+7s+\frac{263}{35}s^2+\frac{121}{63}s^3+\frac{17}{315}s^4$
  & $\frac{787}{45}$ \\
\bottomrule
\end{tabular}
\end{table}

The standard-library verifier implements two independent finite
calculations where such independence is literal: direct adjacent merging for
each boundary permutation and dynamic programming from the split law.  It
also compares the marked tree polynomials with \eqref{eq:markedriccati}, raw
moments with \eqref{eq:momentode}, the Euler-series logarithmic derivative
with the mean coefficients, and the comb count with \eqref{eq:minatom}.
The exact artifact
\path{code/marked_antichain_coefficients.tsv} records $a_n(s)$, the first two
moments, minimum mass, and support size through $n=12$.

These checks do not prove the singularity statements.  Conversely, the
analytic proof does not turn the heap-labeling viewpoint into a second
probability law: both temporal and structural calculations inherit the same
uniform split from \cref{thm:split}.  The honest dependency chain is
\[
 \text{uniform boundaries}
 \Longrightarrow \text{BST split}
 \Longrightarrow
 \begin{cases}
   \text{finite law and moment Riccati hierarchy},\\
   \text{heap-labeled tree and marked antichains}.
 \end{cases}
\]
The two branches cross-check finite coefficients and extremal combs, while
the Sturm--Pringsheim argument is specific to the temporal moment branch.

\section{Ownership boundary and conclusion}

Most importantly, \citet{DisantoEtAl2022} own the same statistic under
$R_n=X_n-1$: the Yule split and finite law, root configurations and unmarked
antichains, the mean Riccati equation and pole, and the second-moment/variance
neighborhood all receive zero credit.  Uniform random-BST splitting,
Cartesian trees, forest ideals, Riccati linearization, Pringsheim's theorem,
and generic singularity extraction are also background.
\citet{AndriantianaWagnerWang2013} own fixed-tree cardinality-marked
antichains, and \citet{ChangFuchs2010,Rosenberg2006} own the caterpillar
probability neighborhood.  The residual recorded here is only the
Yule-averaged bivariate transform and, centrally, the strict pole/radius
continuation for every $r\geq3$.  The all-order identity is a mechanical
interface, not a separate advance.  We do not claim that all equivalent
terminology or unpublished work has been exhausted.

The product-plus-one encoding is useful because it turns the owned
root-configuration statistic into a literal adjacent dynamics.  Retaining a
fixed-tree cardinality marker and averaging it under the Yule law yields the
exact transform, while the triangular Riccati system lets the two low-order
owner results propagate to a strict hierarchy at every higher order.  Those higher moments stop at the exact limsup in
\cref{rem:ceiling}.  Novelty, priority, and external dissemination remain
\textsc{hold} pending specialist review.

\bibliographystyle{plainnat}
\bibliography{references}

\end{document}
